% Unit 110 terjemahan Bahasa Indonesia dari chapter8.tex baris 760--924.
% Sumber: Methods of Algebra, Volume 2, commit
% 9a5803ff77dd3257484cb177f851a73770a59dd3 (CC BY 4.0).
% stable-unit-id: o014.aljabr2.chapter8.dold-kan-a
% Cakupan: Bagian 8.5 dari pembukaan melalui Lema Dold--Kan untuk grup abelian;
% dilanjutkan oleh Unit 111 pada baris 925.

% segment-id: o014.li2.chapter8.dold-kan-a.h001
\section{Korespondensi Dold--Kan}\label{sec:Dold-Kan}
\hypertarget{o014.aljabr2.chapter8.dold-kan-a}{ }
% segment-id: o014.li2.chapter8.dold-kan-a.p001
Mengikuti konvensi topologi aljabar, untuk kategori aditif $\mathcal{A}$ kita
tinjau kategori kompleks rantai di atasnya (Definisi~\ref{def:complex}),
yang ditulis $\cate{Ch}(\mathcal{A})$, beserta subkategorinya
% segment-id: o014.li2.chapter8.dold-kan-a.q001
\[ \cate{Ch}_{\geq 0}(\mathcal{A}) := \left\{ A = (A_n, \partial_n)_{n \in \Z} : \text{kompleks rantai}, \; \forall n < 0, \; A_n = 0 \right\}, \]
dengan $\partial_n=\partial_n^A:A_n\to A_{n-1}$. Untuk objek
$\cate{Ch}_{\geq0}(\mathcal{A})$, jelas cukup memberikan data
$(A_n)_{n\geq0}$ dan $(\partial_n)_{n\geq1}$. Kategori
$\cate{Ch}_{\geq m}(\mathcal{A})$ didefinisikan serupa; bandingkan dengan
versi kompleks kokrantai dalam Definisi~\ref{def:cplx-cat-variant}.
\index[sym1]{ChA@$\cate{Ch}(\mathcal{A})$}

% segment-id: o014.li2.chapter8.dold-kan-a.p002
Dengan memperbesar semesta Grothendieck yang dipilih bila perlu, kita dapat
mengandaikan bahwa $\mathcal{A}$ kategori aditif kecil; lihat
\cite[Asumsi 1.5.2]{Li1}. Dari $\mathcal{A}$ dapat didefinisikan kategori
objek simplisial $\cate{s}\mathcal{A}$; objek simplisial konstan yang
bersesuaian dengan $0$ tetap ditulis
$0\in\Obj(\cate{s}\mathcal{A})$. Bagian ini bertujuan menjelaskan hubungan
dasar antara kategori-kategori tersebut. Lebih tepatnya, kita akan
mendefinisikan tiga funktor
% segment-id: o014.li2.chapter8.dold-kan-a.q002
\[\begin{tikzcd}
	\cate{Ch}_{\geq 0}(\mathcal{A}) \arrow[r, "\Gamma" description] & \cate{s}\mathcal{A} \arrow[bend left, l, "{\mathrm{N}}"] \arrow[bend right, l, "{\mathrm{C}}"']
\end{tikzcd}\]
dengan $\mathrm{N}$ hanya terdefinisi bila $\mathcal{A}$ kategori abelian.
Funktor $\mathrm{C}$ dan $\mathrm{N}$ sama-sama meluas ke objek
semisimplisial dari Catatan~\ref{rem:semisimplicial}; dengan kata lain,
keduanya tidak melibatkan morfisme degenerasi.

% segment-id: o014.li2.chapter8.dold-kan-a.p003
Uraian berikut mengikuti \cite[\S 1.2]{LuHA}. Pertama-tama, kita perjelas
definisi yang berkaitan.

% segment-id: o014.li2.chapter8.dold-kan-a.d001
\begin{definition}[Kompleks rantai tak ternormalisasi]\label{def:unnormalized-chain}
	\index{kompleks rantai tak ternormalisasi@kompleks rantai tak ternormalisasi (unnormalized chain complex)}
	\index{kompleks rantai Moore@kompleks rantai Moore (Moore chain complex)}
	\index[sym1]{CX@$\mathrm{C}X$}
	Untuk objek semisimplisial $X$ dalam $\mathcal{A}$ dan setiap
	$n\in\Z_{\geq1}$, definisikan
	% segment-id: o014.li2.chapter8.dold-kan-a.q003
	\[ \partial_n := \sum_{i=0}^n (-1)^i d_i: X_n \to X_{n-1}. \]
	Objek-objek $(X_n)_{n\geq0}$ beserta $(\partial_n)_n$ membentuk objek
	$\mathrm{C}X$ dari $\cate{Ch}_{\geq0}(\mathcal{A})$, yang disebut
	\emph{kompleks rantai tak ternormalisasi} atau \emph{kompleks rantai
	Moore} yang diberikan oleh $X$. Jadi, $(\mathrm{C}X)_n=X_n$.
\end{definition}

% segment-id: o014.li2.chapter8.dold-kan-a.p004
Perlu ditunjukkan bahwa $\partial_{n-1}\partial_n=0$ bila $n>1$.
Sesungguhnya, identitas
$i<j\implies d_id_j=d_{j-1}d_i$ dalam
\eqref{eqn:simplicial-identity} memberikan
% segment-id: o014.li2.chapter8.dold-kan-a.q004
\begin{multline*}
	\partial_{n-1} \partial_n = \sum_{i=0}^{n-1} \sum_{j=0}^n (-1)^{i+j} d_i d_j \\
	= \sum_{0 \leq i < j \leq n} (-1)^{i+j} d_{j-1} d_i + \sum_{n-1 \geq i \geq j \geq 0} (-1)^{i+j} d_i d_j = 0.
\end{multline*}
Konstruksi ini jelas kanonik dalam $X$ dan memberikan funktor $\mathrm{C}$.

% segment-id: o014.li2.chapter8.dold-kan-a.eg001
\begin{example}\label{eg:classifying-space-std-cplx}
	Misalkan $\Bbbk$ gelanggang komutatif dan
	$\mathcal{A}=\Bbbk\dcate{Mod}$. Untuk grup $\Gamma$, Contoh
	\ref{eg:classifying-space} memberikan himpunan simplisial
	$\mathrm{E}\Gamma$. Definisikan objek simplisial
	$\Bbbk\mathrm{E}\Gamma$ dalam $\Bbbk\dcate{Mod}$ yang suku berderajat
	$n$-nya adalah modul-$\Bbbk$ bebas dengan basis
	$(\mathrm{E}\Gamma)_n$, serta $d_i$ dan $s_i$ diperluas secara linear.
	Akan ditunjukkan bahwa $\mathrm{C}(\Bbbk\mathrm{E}\Gamma)$ persis merupakan
	kompleks rantai $\mathsf{L}:=(\mathsf{L}_n,\partial'_n)_n$ dari
	Definisi~\ref{def:resolution-trivial-module}. Pertama,
	$\mathsf{L}_n$ menurut definisi adalah modul-$\Bbbk$ bebas dengan basis
	$\Gamma^{n+1}$. Ambil pemetaan dari
	$(\mathrm{E}\Gamma)_n=\Gamma^{n+1}$ ke $\mathsf{L}_n$,
	% segment-id: o014.li2.chapter8.dold-kan-a.q005
	\[ \left( g_1, \ldots, g_{n+1} \right) \mapsto \left(g_1 \cdots g_{n+1}, \; \ldots , \; g_n g_{n+1}, \; g_{n+1}\right). \]
	Perluas pemetaan tersebut secara linear pada
	$(\Bbbk\mathrm{E}\Gamma)_n$. Periksa bahwa
	$\sum_i(-1)^id_i$ bersesuaian dengan
	$\partial'_n:\mathsf{L}_n\to\mathsf{L}_{n-1}$, sehingga memberikan
	isomorfisme yang dinyatakan.

	Selain itu, $\Gamma$ beraksi dari kanan pada setiap
	$(\mathrm{E}\Gamma)_n$ melalui
	$(g_1,\ldots,g_{n+1})\xmapsto{g}(g_1,\ldots,g_n,g_{n+1}g)$. Dalam koordinat
	basis $\mathsf{L}_n$, aksi ini menjadi
	$(g_0,\ldots,g_n)\xmapsto{g}(g_0g,\ldots,g_ng)$. Dengan demikian, kita
	memperoleh aksi kanan $\Gamma$ pada $\mathsf{L}$ dari
	Definisi~\ref{def:resolution-trivial-module}. Karena itu, kompleks yang
	diperoleh dengan mengambil koinvarian-$\Gamma$ suku demi suku dari
	kompleks rantai
	$\mathrm{C}(\Bbbk\mathrm{E}\Gamma)$ dan $\mathsf{L}$ juga isomorfik.
	Pada sisi $\mathrm{C}(\Bbbk\mathrm{E}\Gamma)$, operasi ini sama artinya
	dengan mengambil hasil bagi $\mathrm{E}\Gamma$ suku demi suku, lalu
	mengambil modul-$\Bbbk$ bebas. Dengan demikian,
	% segment-id: o014.li2.chapter8.dold-kan-a.q006
	\[ \mathrm{C}(\Bbbk\mathrm{B}\Gamma) \simeq \mathsf{L} \dotimes{\Bbbk[\Gamma]} \Bbbk, \]
	dengan hasil kali tensor melibatkan homomorfisme augmentasi
	$\Bbbk[\Gamma]\to\Bbbk$. Sebagai akibat, homologi grup
	(Definisi~\ref{def:group-coh-ho}) ditafsirkan sebagai
	% segment-id: o014.li2.chapter8.dold-kan-a.q007
	\[ \Hm_n(\mathrm{B}\Gamma; \Bbbk) := \Hm_n \mathrm{C}(\Bbbk\mathrm{B}\Gamma) \simeq \Hm_n(\Gamma, \Bbbk), \quad n \in \Z_{\geq 0}. \]

	Ruas kiri memuat homologi himpunan simplisial, suatu konstruksi umum yang
	akan diperkenalkan dalam Definisi~\ref{def:simplicial-set-homology}.
	Latihan bab ini akan memberikan isomorfisme serupa bagi versi ternormalisasi
	$\overline{\mathsf{L}}$ dari $\mathsf{L}$, yang bersesuaian dengan
	kompleks rantai ternormalisasi
	$\mathrm{N}(\Bbbk\mathrm{E}\Gamma)$ yang akan dibahas dalam
	Definisi~\ref{def:normalized-cplx}.
\end{example}

% segment-id: o014.li2.chapter8.dold-kan-a.d002
\begin{definition}\label{def:DK-Gamma}
	Misalkan $(A_n,\partial_n)_{n\geq0}$ objek dari
	$\cate{Ch}_{\geq0}(\mathcal{A})$. Definisikan objek
	$\Gamma(A)$ dalam $\cate{s}\mathcal{A}$ sebagai berikut:
	% segment-id: o014.li2.chapter8.dold-kan-a.q008
	\[ (\Gamma A)_n := \bigoplus_{t: [n] \twoheadrightarrow [k]} A_k. \]

	Untuk setiap morfisme $f:[m]\to[n]$ dalam $\simpDelta$, morfisme terkait
	$f^*:(\Gamma A)_n\to(\Gamma A)_m$ direpresentasikan, terhadap dekomposisi
	jumlah langsung di atas, oleh matriks
	% segment-id: o014.li2.chapter8.dold-kan-a.q009
	\[ \left( \Phi_{u, t} : A_k \to A_l \right)_{\substack{ u: [m] \twoheadrightarrow [l] \\ t: [n] \twoheadrightarrow [k] }}. \]
	Matriks itu ditentukan sebagai berikut. Untuk suatu
	$t:[n]\twoheadrightarrow[k]$, ambil faktorisasi surjektif--injektif dari
	$tf$,
	% segment-id: o014.li2.chapter8.dold-kan-a.q010
	\begin{equation}\label{eqn:Dold-Kan-epi-mono}\begin{tikzcd}[row sep=tiny]
		& {[n]} \arrow[twoheadrightarrow, rd, "t"] & \\
		{[m]} \arrow[ru, "f"] \arrow[twoheadrightarrow, rd, "u"'] & & {[k]} \\
		& {[l]} \arrow[hookrightarrow, ru, "v"'] &
	\end{tikzcd}\end{equation}
	\begin{itemize}
		\item jika $l=k$, tetapkan $\Phi_{u,t}=\identity$;
		\item jika $l=k-1$ dan $v=\mathrm{d}^0$ (injeksi monoton yang
		melewatkan $0$), tetapkan
		$\Phi_{u,t}=\partial_k:A_k\to A_{k-1}$;
		\item dalam semua kasus lain, tetapkan $\Phi_{u,t}=0$.
	\end{itemize}
	Perhatikan bahwa setelah $f$ dan $t$ diberikan, paling banyak terdapat satu
	$u$ dengan $\Phi_{u,t}\neq0$.
\end{definition}

% segment-id: o014.li2.chapter8.dold-kan-a.p005
Pemeriksaan bahwa $\Gamma A\in\Obj(\cate{s}\mathcal{A})$ agak panjang dan
dihilangkan. Pokok konstruksinya terletak pada suku langsung $A_n$ dari $(\Gamma A)_n$ serta
perilakunya di bawah pemetaan monoton injektif. Semua suku langsung
berderajat lebih rendah $A_k$ berasal dari degenerasi yang diinduksi pemetaan
surjektif $[n]\twoheadrightarrow[k]$; definisi $f^*$ dalam kasus umum hanya
mewujudkan gagasan ini.

% segment-id: o014.li2.chapter8.dold-kan-a.p006
Konstruksi di atas kanonik dan memberikan funktor
$\Gamma:\cate{Ch}_{\geq0}(\mathcal{A})\to\cate{s}\mathcal{A}$.

% segment-id: o014.li2.chapter8.dold-kan-a.d003
\begin{definition}\label{def:normalized-cplx}
	\index{kompleks rantai ternormalisasi@kompleks rantai ternormalisasi (normalized chain complex)}
	\index[sym1]{NX@$\mathrm{N}X$}
	Misalkan $\mathcal{A}$ kategori abelian. Untuk objek semisimplisial $X$ dalam
	$\mathcal{A}$ dan setiap $n\in\Z_{\geq0}$, definisikan
	% segment-id: o014.li2.chapter8.dold-kan-a.q011
	\begin{align*}
		(\mathrm{N} X)_n & := \bigcap_{i=1}^n \Ker\left[ d_i: X_n \to X_{n-1} \right], \quad n \geq 1, \\
		(\mathrm{N} X)_0 & := X_0.
	\end{align*}

	Morfisme $\mathrm{N}X_n\to\mathrm{N}X_{n-1}$ yang diinduksi oleh
	$X_n\xrightarrow{d_0}X_{n-1}$ (periksa menggunakan
	\eqref{eqn:simplicial-identity}), dan selanjutnya ditulis $\partial_n$,
	bersama objek-objek tersebut membentuk kompleks rantai ternormalisasi
	$\mathrm{N}X$ dalam $\cate{Ch}_{\geq0}(\mathcal{A})$.
\end{definition}

% segment-id: o014.li2.chapter8.dold-kan-a.pr001
\begin{proposition}\label{prop:Dold-Kan-u}
	Untuk $X$ seperti di atas dan setiap $n\geq0$, tuliskan
	$u_n:(\mathrm{N}X)_n\hookrightarrow X_n$ untuk inklusi alami. Maka
	$(u_n)_{n\geq0}$ membentuk monomorfisme
	$u:\mathrm{N}X\to\mathrm{C}X$ dalam
	$\cate{Ch}_{\geq0}(\mathcal{A})$; konstruksi ini funktorial dalam $X$.
\end{proposition}
% segment-id: o014.li2.chapter8.dold-kan-a.pf001
\begin{proof}
	Langsung dari definisi $\mathrm{N}X$ dan $\mathrm{C}X$.
\end{proof}

% segment-id: o014.li2.chapter8.dold-kan-a.p007
Beberapa hasil berikut merupakan persiapan bagi Teorema Dold--Kan
\ref{prop:Dold-Kan}.

% segment-id: o014.li2.chapter8.dold-kan-a.le001
\begin{lemma}\label{prop:DK-N-unit}
	Misalkan $\mathcal{A}$ kategori abelian. Untuk setiap
	$A\in\Obj(\cate{Ch}_{\geq0}(\mathcal{A}))$ dan $n\geq0$, berlaku
	$\mathrm{N}\Gamma(A)_n=A_n$. Identifikasi ini memberikan isomorfisme funktor
	$\eta:\identity_{\cate{Ch}_{\geq0}(\mathcal{A})}\rightiso
	\mathrm{N}\Gamma$.
\end{lemma}
% segment-id: o014.li2.chapter8.dold-kan-a.pf002
\begin{proof}
	Jika $t:[n]\twoheadrightarrow[k]$ dan $k<n$, selalu
	dapat dipilih $1\leq i\leq n$ sedemikian sehingga
	$t^{-1}(t(i))$ memiliki sedikitnya dua elemen. Akibatnya,
	$t\mathrm{d}^i$ surjektif dan diagram berikut komutatif:
	% segment-id: o014.li2.chapter8.dold-kan-a.q012
	\[\begin{tikzcd}[row sep=tiny]
		& {[n]} \arrow[twoheadrightarrow, rd, "t"] & \\
		{[n-1]} \arrow[ru, "{\mathrm{d}^i}"] \arrow[twoheadrightarrow, rd, "{t\mathrm{d}^i}"'] & & {[k]} \\
		& {[k]} \arrow[ru, "{\identity}"'] &
	\end{tikzcd}\]
	Definisi~\ref{def:DK-Gamma} kemudian menunjukkan bahwa
	$\mathrm{N}\Gamma(A)_n$ terkandung dalam suku langsung $A_n$ yang
	bersesuaian dengan $\identity:[n]\to[n]$. Definisi tersebut juga segera
	memberikan $\mathrm{N}\Gamma(A)_n=A_n$ dan menunjukkan bahwa
	$\mathrm{N}\Gamma(A)_{n+1}\to\mathrm{N}\Gamma(A)_n$ tepat merupakan
	$\partial_{n+1}:A_{n+1}\to A_n$. Selesailah bukti.
\end{proof}

% segment-id: o014.li2.chapter8.dold-kan-a.le002
\begin{lemma}\label{prop:GammaN-adjunction}
	Misalkan $\mathcal{A}$ kategori abelian. Funktor-funktor di atas membentuk
	adjungsi
	% segment-id: o014.li2.chapter8.dold-kan-a.q013
	\[\begin{tikzcd}
		\Gamma: \cate{Ch}_{\geq 0}(\mathcal{A}) \arrow[shift left, r] & \cate{s}\mathcal{A} : \mathrm{N} \arrow[shift left, l].
	\end{tikzcd}\]
	\begin{itemize}
		\item Morfisme unitnya adalah isomorfisme $\eta$ dari
		Lema~\ref{prop:DK-N-unit}.
		\item Morfisme kounit
		$\epsilon:\Gamma\mathrm{N}\to\identity_{\cate{s}\mathcal{A}}$
		dideskripsikan sebagai berikut. Untuk
		$t:[n]\twoheadrightarrow[k]$, pada suku langsung
		$(\mathrm{N}X)_k$ yang bersesuaian dalam
		$\Gamma\mathrm{N}(X)_n$, pemetaan $\epsilon_{X,n}$ adalah komposisi
		$(\mathrm{N}X)_k\subset X_k\xrightarrow{t^*}X_n$.
	\end{itemize}
\end{lemma}
% segment-id: o014.li2.chapter8.dold-kan-a.pf003
\begin{proof}
	Untuk $A\in\Obj(\cate{Ch}_{\geq0}(\mathcal{A}))$ dan
	$X\in\Obj(\cate{s}\mathcal{A})$, pertama akan dibuktikan bahwa komposisi
	% segment-id: o014.li2.chapter8.dold-kan-a.q014
	\[ \Hom_{\cate{s}\mathcal{A}}(\Gamma(A), X) \xrightarrow{\mathrm{N}} \Hom_{\cate{Ch}_{\geq 0}(\mathcal{A})}(\mathrm{N}\Gamma(A), \mathrm{N}(X)) \xrightarrow{\eta_A^*} \Hom_{\cate{Ch}_{\geq 0}(\mathcal{A})}(A, \mathrm{N}(X)) \]
	merupakan bijeksi. Inversnya didefinisikan secara konkret sebagai berikut. Untuk
	$\phi=(\phi_n)_{n\geq0}:A\to\mathrm{N}(X)$, definisikan morfisme
	% segment-id: o014.li2.chapter8.dold-kan-a.q015
	\[ \Phi_n: (\Gamma A)_n = \bigoplus_{t: [n] \twoheadrightarrow [k]} A_k \to X_n \]
	dengan pembatasan pada suku langsung yang bersesuaian dengan
	$t:[n]\twoheadrightarrow[k]$ berupa komposisi
	% segment-id: o014.li2.chapter8.dold-kan-a.q016
	\[ A_k \xrightarrow{\phi_k} (\mathrm{N}X)_k \subset X_k \xrightarrow{t^*} X_n. \]
	Untuk $t$ seperti di atas dan sebarang $f:[m]\to[n]$, ambil faktorisasi
	surjektif--injektif dari $tf$ seperti dalam
	\eqref{eqn:Dold-Kan-epi-mono}, lalu tinjau diagram dalam $\mathcal{A}$
	(lihat Definisi~\ref{def:DK-Gamma}):
	% segment-id: o014.li2.chapter8.dold-kan-a.q017
	\[\begin{tikzcd}[column sep=large]
		A_k \arrow[r, "\phi_k"] \arrow[d, "{\Phi_{u, t}}"'] & (\mathrm{N}X)_k \arrow[hookrightarrow, r] \arrow[d, "{v^*|_{(\mathrm{N}X)_k}}"] & X_k \arrow[r, "{t^*}"] \arrow[d, "{v^*}"] & X_n \arrow[d, "{f^*}"] \\
		A_l \arrow[r, "\phi_l"'] & (\mathrm{N}X)_l \arrow[hookrightarrow, r] & X_l \arrow[r, "{u^*}"'] & X_m.
	\end{tikzcd}\]
	Persegi kiri komutatif menurut definisi kompleks rantai ternormalisasi
	$\mathrm{N}X$, persegi tengah jelas komutatif, dan persegi kanan komutatif
	karena $tf=vu$. Jadi, $\Phi=(\Phi_n)_n$ adalah morfisme dalam
	$\cate{s}\mathcal{A}$.

	Verifikasi bahwa $\eta_A^*\mathrm{N}$ dan $\phi\mapsto\Phi$ saling invers
	bersifat langsung. Dalam deskripsi $\phi\mapsto\Phi$, ambil
	$\phi=\identity_{\mathrm{N}X}$; hasilnya tepat $\epsilon$ yang dinyatakan.
\end{proof}

% segment-id: o014.li2.chapter8.dold-kan-a.le003
\begin{lemma}\label{prop:DK-Ab}
	Ambil $\mathcal{A}=\cate{Ab}$. Maka $\Gamma$ merupakan ekuivalensi kategori; lebih
	tepatnya, pasangan adjoin
	% segment-id: o014.li2.chapter8.dold-kan-a.q018
	$\begin{tikzcd}
		\Gamma: \cate{Ch}_{\geq 0}(\cate{Ab}) \arrow[shift left, r] & \cate{s}\cate{Ab} : \mathrm{N} \arrow[shift left, l]
	\end{tikzcd}$
	merupakan adjungsi yang sekaligus merupakan ekuivalensi dalam pengertian
	\cite[Teorema 2.6.12]{Li1}.
\end{lemma}
% segment-id: o014.li2.chapter8.dold-kan-a.pf004
\begin{proof}
	Cukup dibuktikan bahwa morfisme kounit $\epsilon$ dari
	Lema~\ref{prop:GammaN-adjunction} adalah isomorfisme. Misalkan
	$X\in\Obj(\cate{s}\mathcal{A})$. Akan dibuktikan bahwa
	$\epsilon_{X,n}:\Gamma(\mathrm{N}X)_n\to X_n$ injektif untuk setiap
	$n\geq0$. Misalkan $x=(x_t)_t$ elemen domainnya, dengan $t$ merentang
	semua pemetaan monoton surjektif $[n]\twoheadrightarrow[k]$. Untuk suatu
	$t$, definisikan injeksi monoton
	% segment-id: o014.li2.chapter8.dold-kan-a.q019
	\begin{align*}
		s = s_t : [k] & \hookrightarrow [n] \\
		i & \mapsto \min t^{-1}(i),
	\end{align*}
	yang memenuhi $ts=\identity_{[k]}$. Andaikan $x\neq0$, dan tuliskan
	% segment-id: o014.li2.chapter8.dold-kan-a.q020
	\[ S := \left\{t: x_t \neq 0 \right\} \neq \emptyset. \]
	Pilih $k$ terkecil sehingga terdapat $t:[n]\twoheadrightarrow[k]$ dalam
	$S$, kemudian pilih $t$ seperti itu yang meminimalkan
	$\sum_{i=0}^k\min t^{-1}(i)$, dan konstruksikan $s$. Kita akan membuktikan
	bahwa $s^*(\epsilon_{X,n}(x))=x_t\in X_k$, sehingga
	$\epsilon_{X,n}(x)\neq0$.

	Menurut deskripsi konkret $\epsilon_{X,n}$, cukup ditunjukkan bahwa untuk
	setiap $t':[n]\twoheadrightarrow[k']$,
	$s^*(t')^*(x_{t'})\neq0$ mengakibatkan $t=t'$. Tetapkan
	$u:=t's:[k]\to[k']$. Jika $t'\notin S$, maka $x_{t'}=0$. Maka cukup
	meninjau kasus $t'\in S$ dengan
	$s^*(t')^*(x_{t'})=u^*(x_{t'})\neq0$.

	Karena $t$ monoton dan surjektif, $\min t^{-1}(0)=0$, dan demikian pula
	untuk $t'$, sehingga $u(0)=0$. Selain itu,
	$x_{t'}\in(\mathrm{N}X)_{k'}$ mengakibatkan
	$u^*(x_{t'})\neq0$ hanya mungkin bila
	$\Image(u)\supset\{1,\ldots,k'\}$. Jadi, kita dapat mengandaikan $u$
	surjektif. Minimalitas $k$ kemudian memaksa $k'=k$ dan
	$u=\identity$.

	Kesimpulannya,
	$t'(\min t^{-1}(i))=i$ untuk $i=0,\ldots,k$, sehingga
	$\min(t')^{-1}(i)\leq\min t^{-1}(i)$. Dari pilihan $t$, diperoleh
	$\min(t')^{-1}(i)=\min t^{-1}(i)$ untuk setiap $0\leq i\leq k$.
	Karena pemetaan monoton surjektif ditentukan oleh titik awal
	serat-seratnya, ini sama artinya dengan $t=t'$. Jadi,
	$\epsilon_{X,n}$ injektif.

	Selanjutnya kita buktikan secara induktif pada $n$ bahwa $\epsilon_{X,n}$
	surjektif. Akan ditunjukkan bahwa untuk setiap $0\leq i\leq n$,
	% segment-id: o014.li2.chapter8.dold-kan-a.q021
	\[ \Image\left( \epsilon_{X, n} \right) \supset X(i)_n := \bigcap_{i < j \leq n} \Ker(d_j) \subset X_n. \]

	Untuk $i=0$, berlaku $X(i)_n=(\mathrm{N}X)_n$, dan inklusi di atas mudah
	diperoleh dari deskripsi konkret $\epsilon_{X,n}$. Tujuan kita ialah kasus
	$i=n$. Misalkan $n\geq i\geq1$ dan $y\in X(i)_n$. Hipotesis induksi untuk
	$n-1$ dan deskripsi $\epsilon_{X,n}$ memberikan
	% segment-id: o014.li2.chapter8.dold-kan-a.q022
	\[ s_{i-1} d_i(y) \in s_{i-1}\left( \Image\left(\epsilon_{X, n-1}\right) \right) \subset \Image\left( \epsilon_{X, n} \right). \]

	Di sisi lain, identitas simplisial \eqref{eqn:simplicial-identity}
	memberikan
	% segment-id: o014.li2.chapter8.dold-kan-a.q023
	\begin{gather*}
		d_i s_{i-1} d_i = d_i, \\
		j > i \implies d_j s_{i-1} d_i = s_{i-1} d_{j-1} d_i = s_{i-1} d_i d_j.
	\end{gather*}
	Jadi, $y-s_{i-1}d_i(y)\in X(i-1)_n$, dan hipotesis induksi untuk $i-1$
	menunjukkan bahwa elemen ini berada dalam
	$\Image(\epsilon_{X,n})$. Dengan demikian,
	$y\in\Image(\epsilon_{X,n})$. Selesailah bukti.
\end{proof}
